\section{A division-free integral basis and a unit relation minor}
\label{ir:sec:integral}

\subsection{Choosing the surviving grid points}
Write $U_m=(\Z/m\Z)^\times$ and use a one-element set $U_1$.
For every prime $p$ and exponent $a\ge1$, choose a distinguished lift
in each fiber of $U_{p^a}\to U_{p^{a-1}}$. At $a=1$, this means choosing
one member of $U_p$. Let $E_{p,a}$ be the remaining members. Thus
\begin{equation}\label{ir:eq:local-count}
 |E_{p,a}|=\varphi(p^a)-\varphi(p^{a-1}),\qquad \varphi(1)=1.
\end{equation}
The following fixed choice makes the construction executable:
\begin{align}
 E_{p,1}&=\{2,3,\ldots,p-1\},\label{ir:eq:local-choice}\\
 E_{p,a}&=\{b+jp^{a-1}:1\le b<p^{a-1},\ p\nmid b,
                         \ 1\le j\le p-1\}\quad(a\ge2).\nonumber
\end{align}
For $p=2$, the first set is empty. This is intentional.

Every point $x$ of exact denominator $m$ has a unique additive primary
expansion
\begin{equation}\label{ir:eq:primary-coordinates}
 x=\sum_{p^a\parallel m}\frac{b_p}{p^a}\pmod\Z,
 \qquad b_p\in U_{p^a}.
\end{equation}
For $x=n/m$ in lowest terms,
$b_p\equiv n(m/p^a)^{-1}\pmod{p^a}$. Define $\B_q$ to contain $0$
and all points of denominator $m\mid q$ for which
$b_p\in E_{p,a}$ at every nonzero primary component. We call a primary
component \emph{bad} when its residue is the distinguished, discarded
lift. This adjective refers only to a basis choice.

The cardinality telescopes prime by prime:
\begin{equation}\label{ir:eq:basis-count}
 |\B_q|=
 \prod_{p^e\parallel q}\left(1+
   \sum_{a=1}^{e}\bigl(\varphi(p^a)-\varphi(p^{a-1})\bigr)\right)
 =\varphi(q).
\end{equation}
The point $0$ must be retained for this count and for the normal forms.

\begin{theorem}[Integral normal form]\label{ir:thm:integral}
The images of $e_b$, $b\in\B_q$, form an $R_q$-basis of $\D_q$.
Consequently, for every commutative ring $S$ and every choice of weights
$a_p\in S$, the same points form an $S$-basis of $\D_q(S,a)$.
In particular this module is free of rank $\varphi(q)$ even in positive
characteristic, at zero weights, and over nonreduced rings.
\end{theorem}
The spanning algorithm is immediate and is proved next. Its independence
is established in Section~\ref{ir:sec:resolution}, without dividing by
any group order or weight.

\subsection{The terminating straightening rule}
Let $y\notin\B_q$, let $m=\den(y)$, and choose any bad prime $p$ for
its primary coordinates. The row with target $py$ gives
\begin{equation}\label{ir:eq:straighten}
 e_y=A_pe_{py}-\sum_{\substack{pz=py\\z\ne y}}e_z
       \quad\text{in }\D_q.
\end{equation}
All terms on the right have denominator dividing $m$.
At exact denominator $m$, the $p$-primary components of the other
solutions are the allowed lifts in the same fiber; all other primary
components are unchanged. When the $p$-component has order $p$, there
is also one solution with zero $p$-component, hence lower denominator.
The point $py$ itself has denominator $m/p$.

Order points first by their positive denominator and then by the number
of bad primary components. Every right-hand term of
\eqref{ir:eq:straighten} is strictly earlier in this order. Recursion
therefore terminates, and proves that $\B_q$ spans. The implementation
always chooses the smallest bad prime, but the theorem permits any
choice. Independence will imply that all choices yield the same
normal form.

\begin{proposition}[Support and polynomial degree]\label{ir:prop:support}
Write the normal form of $e_y$, with $m=\den(y)$, as
\[
 \NF(e_y)=\sum_{b\in\B_q}c_{y,b}(\mathbf A)e_b.
\]
A nonzero coefficient has $\den(b)=d\mid m$, and
\begin{equation}\label{ir:eq:degree-bound}
 \deg_{A_p}c_{y,b}\le v_p(m/d)\qquad(p\mid q).
\end{equation}
All coefficients belong to $\Z[\mathbf A]$.
\end{proposition}
\begin{proof}
The recursion never increases a primary denominator. Multiplication by
$A_p$ occurs only in the term $e_{py}$, where that denominator has
lost one power of $p$. Every path from denominator $m$ to denominator
$d$ can therefore contain at most $v_p(m/d)$ such multiplications.
The other coefficients in the recursion are $-1$. Summing the finitely
many paths proves the assertions, including possible cancellations.
\end{proof}

There is no assertion here that the coefficient of $e_0$ factors as a
product of $A_p-1$. At level $15$ the coefficient in one normal form is
$1-A_3A_5$; this distinction will matter at the spectral pole.

\subsection{An explicit maximal minor equal to one}
For each nonbasis point $y$, select the single row $r_{p,py}$ used in
its first straightening step. Select the nonbasis points as columns,
and order rows by their selected pivot. The restricted matrix is
triangular, with every diagonal entry equal to $1$: all other nonbasis
columns are earlier in the terminating order.

The selected rows are distinct. Indeed a selected row has its chosen
point as its unique latest nonbasis point, with coefficient $1$. Two
identical rows could not have different such latest points.
We have proved the following stronger integrality statement.

\begin{theorem}[Universal unit minor and relation basis]
\label{ir:thm:minor}
The prime-relation matrix contains an explicitly selected
$(q-\varphi(q))\times(q-\varphi(q))$ minor equal to $1$ in $R_q$.
The selected rows form an $R_q$-basis of the full relation module
$\R_q$. After specializing the weights to integers, all nonzero Smith
invariants of the complete prime-relation matrix are $1$; there are
exactly $q-\varphi(q)$ of them.
\end{theorem}
\begin{proof}
Triangularity proves the minor statement. Selected-row elimination
reduces every vector to a combination of basis-point symbols. By
Theorem~\ref{ir:thm:integral}, a vector in $\R_q$ has zero such
remainder, so every raw row is a combination of the selected ones.
The minor proves their independence. The same argument specializes
without division. Equivalently, after an integer specialization the
image lattice is saturated and has the indicated rank; the unit minor
also proves directly that all Smith invariants are units.
\end{proof}
The proof order is not circular: the resolution below proves
Theorem~\ref{ir:thm:integral} independently of the assertion that the
selected rows span all relations.

\subsection{Certificates rather than opaque elimination}
The straightener can retain the exact row combination at every step.
If $C_y$ denotes a certificate for $e_y-\NF(e_y)$, then
\begin{equation}\label{ir:eq:certificate-recursion}
 C_y=r_{p,py}+A_pC_{py}-\sum_{\substack{pz=py\\z\ne y}}C_z,
 \qquad C_b=0\ (b\in\B_q).
\end{equation}
A verifier need not trust the recursion. It can expand the listed raw
rows in the original $q$-generator free module, multiply their explicit
integer polynomials, and compare coefficients. The accompanying
verifier does exactly this for the three frozen identities.

\begin{example}[Small levels]
The specified basis choices give
\[
 \B_4=\{0,3/4\},\qquad
 \B_{12}=\{0,5/12,2/3,3/4\},
\]
\[
 \B_{15}=\{0,1/15,4/15,2/5,7/15,3/5,2/3,4/5\}.
\]
Also $\B_{30}=\B_{15}$ as sets of rational points. An exact primary
component of order two has no allowed residue, which explains this
last equality. It does not make the weight $A_2$ irrelevant to all
normal-form coefficients.
\end{example}

The procedure terminates with polynomial coefficients, but we do not
assert a polynomial-time bound in the bit length of $q$. There are
$q$ raw grid points to begin with, and certificate coefficients can
grow. The finite support and degree bounds are rigorous; sharper
complexity and coefficient-height estimates remain separate questions.
